\section{Worked Example: How a Heavy Chain Emerges}
Suppose one frontier level emits successors whose discrete keys are:
\[
\underbrace{dk_a,dk_a,dk_a}_{3},\underbrace{dk_b,dk_b}_{2},\underbrace{dk_c,\dots,dk_c}_{40},\underbrace{dk_d}_{1},\underbrace{dk_e,\dots,dk_e}_{18}.
\]
Then pending counts are
\[
|\Pend_{dk_a}|=3,\ |\Pend_{dk_b}|=2,\ |\Pend_{dk_c}|=40,\ |\Pend_{dk_d}|=1,\ |\Pend_{dk_e}|=18.
\]
If prior live sizes are
\[
|\Live_{dk_a}|=2,\ |\Live_{dk_b}|=1,\ |\Live_{dk_c}|=55,\ |\Live_{dk_d}|=1,\ |\Live_{dk_e}|=12,
\]
a useful first-order cost proxy is
\[
\cost(dk) = \alpha\,|\Pend_{dk}|\,|\Live_{dk}| + \beta\,|\Pend_{dk}|^2.
\]
So
\[
\cost(dk_c) \gg \cost(dk_a),\cost(dk_b),\cost(dk_d),\cost(dk_e).
\]
Even with many threads, one worker processing $dk_c$ can dominate level completion.

\paragraph{How to read this example.}
Keys $dk_a,dk_b,\dots$ are not equally hard.
The chain for $dk_c$ is hard for two reasons at once:
\begin{enumerate}
	\item many new candidates arrive this level ($|\Pend_{dk_c}|=40$),
	\item many old survivors already exist ($|\Live_{dk_c}|=55$).
\end{enumerate}
So each new candidate may scan a large live set, and then possibly trigger merge work.
This is the mechanism behind ``one hot chain slows the whole level.''

\paragraph{Practical takeaway.}
If your profiler shows a long tail at level barriers, inspect the top few $dk$ by pending and live sizes.
In many runs, a tiny number of keys account for most settle time.

\begin{figure}[t]
\centering
\begin{tikzpicture}[>=Latex, font=\small]
  \tikzstyle{box}=[draw, rounded corners, minimum width=33mm, minimum height=7mm]
  \node[box] (p0) at (0,0) {Pending for $dk_c$: 40 zones};
  \node[box] (l0) at (0,-1.2) {Live antichain: 55 zones};
  \node[box] (m0) at (0,-2.4) {Settle: subsume/merge scans};
  \node[box] (o0) at (0,-3.6) {Survivors feed next frontier};

  \draw[->] (p0) -- (m0);
  \draw[->] (l0) -- (m0);
  \draw[->] (m0) -- (o0);

  \node at (5.6,-1.8) {$\cost(dk_c) \approx \alpha p\,l + \beta p^2$};
\end{tikzpicture}
\caption{A heavy chain forms when both pending and live sets are large.}
\label{fig:chain-growth}
\end{figure}

